\begin{theorem}[Unconditional large-order missing-color bound on the current ladder]
Let $G=(V,E)$ be a finite loopless multigraph, let $e_0\in E$, and let
$k$ satisfy $\Delta(G)+1\leq k$.  Let $T$ be an arbitrary-input
positive-rung ETT with rung number $n$, generating coloring $\varphi_n$,
and certified Tashkinov series, and suppose that $T$ has maximum property.
Assume a technical induction context through rung $n$ and assume
$11\leq t(G)$, where $t(G)$ is the Tashkinov order.  If $T_n$ denotes
the exact terminal ladder tree
\[
 T_n=T.\mathrm{series.data.tree}
       ((\operatorname{Fin.last}n).\mathrm{castSucc}),
\]
then
\[
 2n+11\leq |\overline{\varphi_n}(T_n)|.
\]
The assertion is uniform in the certified input coloring and assumes no
hierarchy base, good hierarchy, or prior missing-color estimate.
\end{theorem}
\begin{proof}
For $1\leq i\leq n$, write
\[
 T_i=T.\mathrm{series.data.tree}(i),\qquad
 \varphi_{i-1}=T.\mathrm{series.data.coloring}(i-1),
\]
using the canonical finite indices.  The positivity clause of
\noderef{ExtendedTashkinovTree} ensures $n\geq1$.

We first prove that, for every $1\leq i\leq n$,
\[
 V(T_i)\text{ is elementary under }\varphi_{i-1}.             \tag{1}
\]
For $i=1$, this is exactly the first conclusion of
\noderef{TechnicalTheoremBase} applied to the certified series.  Now let
$2\leq i\leq n$ and set $r=i-1$.  Then $0<r<n$.  Apply
\noderef{CertifiedETTRestriction} to the given maximum ETT with target
rung $r$.  Its statement supplies an actual positive-rung ETT with the
same arbitrary input coloring, rung exactly $r$, displayed tree exactly
the original $T_{r+1}=T_i$, generating coloring exactly the original
$\varphi_r=\varphi_{i-1}$, and inherited maximum property.  The
\texttt{lowerOne} field of \noderef{TechnicalInductionContext}, applied
to this certified lower ETT and $0<r<n$, gives (1).

Next inspect the exact certificate at every positive data slot
$i\in\{1,\ldots,n\}$ in \noderef{TashkinovSeries}.  The initial
constructor of \noderef{TashkinovSeriesStep} cannot occur at a positive
slot.  Each of the SE, PE, and RE constructors has the following common
fields: $\varphi_{i-1}$ is proper and leaves exactly $e_0$ uncolored,
$V(T_i)$ is closed under $\varphi_{i-1}$, $T_i$ is an exact ordered
initial segment of $T_{i+1}$, and $T_i$ has strictly smaller sequence
length than $T_{i+1}$.  The unique-uncolored field in particular gives
$\varphi_{i-1}(e_0)$ uncolored.

Apply \noderef{TreeMissingColorsCardLowerBound} to $T_i$, using (1),
$\varphi_{i-1}(e_0)$ uncolored, and $\Delta(G)+1\leq k$.  It gives
\[
 |V(T_i)|+2\leq|\overline{\varphi_{i-1}}(T_i)|,               \tag{2}
\]
so choose
$\alpha_i\in\overline{\varphi_{i-1}}(T_i)$.  We claim that $T_i$ is
$\{\alpha_i\}$-closed.  Indeed, choose a vertex $v\in V(T_i)$ missing
$\alpha_i$.  If a boundary edge identity $f$ had color $\alpha_i$, the
\noderef{ColoredVertexSetClosed} field just extracted, applied to $f$ and
$v$, would say that $\alpha_i$ is not missing at $v$, a contradiction.
Thus the $\alpha_i$-colored boundary set is empty, which is precisely
\noderef{TreeColorSetClosed} for the singleton $\{\alpha_i\}$.  Now
\noderef{ElementaryClosedMissingColorOdd}, using properness, (1), the
choice of $\alpha_i$, and this singleton closedness, yields
\[
                         |V(T_i)|\text{ is odd}.               \tag{3}
\]

For $1\leq i<n$, the ordered-extension and strict-length fields of the
same slot-$i$ certificate, followed by
\noderef{TreeSequenceInitialSegmentVertexSetStrict}, give
$V(T_i)\subsetneq V(T_{i+1})$.  Hence
$|V(T_i)|<|V(T_{i+1})|$.  Both cardinalities are odd by (3).  Writing
them as $2a+1$ and $2b+1$, strict inequality gives $a<b$, and therefore
\[
                         |V(T_i)|+2\leq|V(T_{i+1})|.            \tag{4}
\]
The first clause of \noderef{MaximumProperty} says
$|V(T_1)|=t(G)$, so the large-order hypothesis gives $11\leq|V(T_1)|$.
Induction on $i$, with this base and recurrence (4), proves
\[
                         2i+9\leq|V(T_i)|
       \quad(1\leq i\leq n).                                  \tag{5}
\]
In particular, $2n+9\leq|V(T_n)|$.

It remains to use the generating coloring $\varphi_n$, rather than the
preceding coloring $\varphi_{n-1}$ used in (1)--(5).  Let $a$ be the last
positive data index and inspect the exact certificate at $a$.  Again it
cannot be initial.  The generating-color equality in
\noderef{ExtendedTashkinovTree} identifies $\varphi_n$ with the current
coloring stored at $a$.  We prove
\[
 V(T_n)\text{ is elementary under }\varphi_n,                 \tag{6}
\]
and simultaneously record that $\varphi_n$ is proper and leaves $e_0$
uncolored, by considering its three possible constructors.

If the last constructor is SE, its
\noderef{MaximumDefectiveChoice} has a selected coloring which is proper,
leaves exactly $e_0$ uncolored, and is stable on $T_n$ relative to
$\varphi_{n-1}$.  The SE current-coloring equation identifies this
selected coloring with $\varphi_n$.  The vertex clause of
\noderef{StableColoring} identifies, at every vertex of $T_n$, its
missing-color set with the one under $\varphi_{n-1}$.  Thus (1) for
$i=n$ transfers directly to (6), and the selected-coloring fields supply
properness and the root condition.

If the last constructor is RE, its current-coloring equation is
$\varphi_n=\varphi_{n-1}$.  Thus (6), properness, and the root condition
are exactly (1) and the common preceding-coloring fields already
extracted from that RE constructor.

Suppose finally that the last constructor is PE.  Its current fields say
directly that $\varphi_n$ is proper and leaves exactly $e_0$ uncolored.
By \noderef{HierarchyBaseExists}, choose a hierarchy base for this ETT
under the generating coloring $\varphi_n$.  Assemble the six lower-rung
premises of \noderef{ParallelClosureUnionElementary} point by point from
the fields \texttt{lowerOne}, \texttt{lowerTwo}, \texttt{lowerThree},
\texttt{lowerFour}, \texttt{lowerFive}, and \texttt{lowerSix} of the
given \noderef{TechnicalInductionContext}; their quantifiers all permit
arbitrary certified input colorings and their common hypotheses are
exactly positive rung, rung below $n$, and maximum property.  Apply
\noderef{ParallelClosureUnionElementary} with $\sigma=\varphi_n$.  The
required stability is reflexive by \noderef{sc2}; the PE-kind equation,
properness, root condition, and positive rung have just been supplied.
It follows that the ordered-union core of the chosen base is elementary
under $\varphi_n$.

The PE clause of \noderef{HierarchyBase} says that this core is the exact
ordered union of $T_n$ with the support-prefix closure.  In particular,
$T_n$ is an ordered initial segment of the core, so
\noderef{TreeSequenceInitialSegmentVertexSetSubset} gives
$V(T_n)\subseteq V(\text{core})$.  Restricting the pairwise-disjoint
missing-color condition in \noderef{ColoredVertexSetElementary} from the
core to this subset proves (6).

All three cases therefore give (6), properness, and in particular
$\varphi_n(e_0)$ uncolored.  Apply
\noderef{TreeMissingColorsCardLowerBound} once more, now to $T_n$ under
$\varphi_n$.  We obtain
\[
 |V(T_n)|+2\leq|\overline{\varphi_n}(T_n)|.                   \tag{7}
\]
Combining (5) at $i=n$ with (7) gives
\[
 2n+11=(2n+9)+2
   \leq |V(T_n)|+2
   \leq |\overline{\varphi_n}(T_n)|,
\]
which is the desired bound.
\end{proof}
